\section{Síntesis y ampliación}

\subsection{Tabla de síntesis}

\begin{table}[H]
\centering
\begin{tabular}{lp{8cm}p{4cm}}
\toprule
Sección & Enfoque & Recomendación posterior \\
\midrule
Repaso de los fundamentos del MDP & Relación entre la iteración de Bellman y la iteración de política/value & Revisar de nuevo el teorema de contracción y sus formas de evolución \\
Q-Learning tabular & Convergencia off-policy y equilibrio de la exploración & Revisar la cobertura del replay buffer en las tareas actuales \\
Fitted Q-Iteration y DQN & Aproximación de funciones, técnicas de estabilización y detalles de optimización & Verificar en sistemas existentes el efecto de target network / double DQN \\
De DQN a Actor-Critic & Transición hacia continuous control y la arquitectura actor-critic & Comparar la estimación del gradiente de deterministic policy gradient con la de Q-learning \\
Práctica de sistemas y monitoreo & Métricas, estado del replay buffer, mecanismos de intervención & Establecer un mecanismo de instrumentation revoke \\
Variantes y costos & Rainbow, modelo de costos, pila de inferencia & Planificar el presupuesto de entrenamiento/inferencia y conservar una ruta de reincorporación \\
Casos de aplicación y gobernanza & Avances en benchmark y sim-to-real, aspectos de gobernanza & Planificar soft/hard interrupt y la reversión ante anomaly \\
\bottomrule
\end{tabular}
\caption{Contenidos principales de esta clase y puntos de entrada para la ingeniería}
\end{table}

\subsection{Lecturas para profundizar}

\begin{knowledgebox}{Material de ampliación}
\begin{itemize}
    \item Watkins \& Dayan, ``Q-Learning,'' Machine Learning 1992
    \item Mnih et al., ``Human-level Control through Deep Reinforcement Learning (DQN),'' Nature 2015
    \item Van Hasselt et al., ``Deep Reinforcement Learning with Double Q-Learning,'' AAAI 2016
    \item Lillicrap et al., ``Continuous Control with Deep Reinforcement Learning (DDPG),'' ICLR 2016
    \item Sutton \& Barto, ``Reinforcement Learning: An Introduction,'' 2nd Edition, 2018
    \item Documentación de OpenAI Safety Gym 2.0 (permite hacer pruebas de interruptibility)
\end{itemize}
\end{knowledgebox}

\subsection{Resumen de la sección}

Esta clase parte del MDP y la ecuación de Bellman para desarrollar, paso a paso, la evolución de Q-Learning: de la forma tabular a las redes profundas y, después, a Actor-Critic. A lo largo del recorrido intercala la práctica de sistemas, las variantes y su COST, así como casos de gobernanza y de aplicación, y ofrece un checklist completo de repaso y de operación.

\end{document}
